\chapter{Conclusion et perspectives}
\label{chap:conclusion-perspectives}
\justifying
\sloppy

\section{Synthèse technique}

Le projet \textbf{SPADMIC} est structuré comme un prototype ASIC de lecture pour matrice de \textbf{SPAD}, dans lequel la mesure temporelle constitue une brique essentielle de la reconstruction. Le besoin système impose plusieurs axes de lecture, un flux de sortie compact et une observabilité suffisante pour permettre une calibration externe côté hôte.

Le choix d'une architecture Vernier multiphase répond à une contrainte centrale : obtenir une finesse temporelle élevée sans supposer qu'un délai élémentaire unique suffise à porter toute la résolution. L'observation parallèle de plusieurs phases réduit le temps nécessaire à la décision de mesure et prépare naturellement une lecture sous forme de matrice de phase.

L'implémentation \textbf{RTL} active transforme ce principe en bloc intégrable. Le cœur \textbf{MPTDC} associe une capture asynchrone des événements, deux familles de phases Vernier, une matrice de détection, un snapshot de l'image de mesure, une banque à deux contextes, une lecture en domaine système et une trame de sortie compacte sur \SI{16}{\bit}. L'objectif du bloc n'est donc pas seulement de produire un temps brut : il est de fournir, pour les trois axes de lecture \textbf{SPADMIC}, une mesure temporelle exploitable et compatible avec une correction hors ASIC.

La vérification fonctionnelle a consolidé cette architecture. Le bloc a été exercé avec une démarche proche de l'industrie : bancs unitaires, bancs d'intégration, scénarios dirigés et aléatoires, comparaison systématique avec un modèle de référence, contrôle de la cohérence des paquets, de la rétropression, des resets, des timeouts et des cas limites. Cette étape a permis de vérifier le contrat numérique du bloc avant de s'appuyer sur ses sorties pour la caractérisation.

\section{Résultats principaux}

La caractérisation \textbf{RTL}/\textit{Xcelium} confirme que la structure brute du Vernier multiphase ne doit pas être interprétée comme une échelle temporelle directement uniforme. Avant calibration, l'erreur présente un biais important, avec une RMSE autour de \SI{1.94}{\nano\second}. Ce comportement est cohérent avec la non-uniformité de la grille fine et justifie le choix de transmettre une trame suffisamment riche pour une correction externe.

Après application de la LUT hôte, la validation tenue hors apprentissage atteint une RMSE de \SI{18.56}{\pico\second}, avec un percentile 99 de \SI{38.99}{\pico\second}. Ces valeurs ne constituent pas une mesure silicium, mais elles montrent qu'en simulation RTL les informations exportées par le protocole compact suffisent à supprimer le biais principal et à recentrer l'erreur autour de zéro. Le moyennage statistique peut ensuite améliorer la précision dans un domaine réaliste, notamment autour des seuils de \SI{10}{\pico\second} et \SI{5}{\pico\second}, sans que cela soit présenté comme une promesse de performance à très grand nombre d'échantillons.

\section{Perspectives}

La portée des résultats doit rester sobrement délimitée. Le bloc atteint un état cohérent de fin de conception numérique pré-silicium : l'architecture est structurée, le protocole est stabilisé, la vérification fonctionnelle est avancée et la calibration externe est démontrée en simulation. Les étapes structurantes restantes concernent la validation physique, la fermeture temporelle et, à terme, la confrontation aux mesures silicium.

Les travaux en cours portent donc sur l'implémentation physique et le \textit{timing closure}. Cette phase ne consiste pas seulement à lancer une synthèse ; elle implique une réorganisation progressive de la profondeur logique, la réduction des chemins critiques, la correction des violations \textit{setup}/\textit{hold}, la maîtrise du \textit{skew}, l'affinage des contraintes de synthèse et de \textbf{STA}, puis l'adaptation du \textbf{RTL} aux réalités du placement-routage (\textit{Place and Route}, \textbf{PnR}). Ces ajustements permettront de vérifier si les choix numériques retenus restent robustes une fois soumis aux contraintes de placement, de routage et d'horlogerie.

La suite devra également consolider les modèles d'oscillateurs, rejouer la caractérisation après synthèse puis après layout, formaliser le flot de génération des LUT et préciser les règles de rejet ou de repli pour les paquets hors domaine de calibration. La comparaison avec les mesures silicium restera l'étape décisive pour transformer les résultats \textbf{RTL}/\textit{Xcelium} en performance expérimentale.

\section{Apport méthodologique}

Au-delà du fonctionnement du bloc, la démarche suivie couvre une chaîne complète d'ingénierie numérique ASIC. Les choix d'architecture, la micro-architecture, l'écriture \textbf{RTL}, la vérification, la caractérisation, le timing et l'intégration physique ne sont pas des étapes indépendantes : chacune impose des contraintes qui peuvent remettre en question les hypothèses prises aux étapes précédentes.

L'expérience montre ainsi que concevoir un bloc ne revient pas seulement à obtenir une simulation fonctionnelle correcte. Il faut raisonner en boucle entre le besoin système, la structure matérielle, les observables disponibles, la vérifiabilité, la profondeur logique, les chemins critiques et les contraintes d'intégration. Cette approche donne une base plus solide pour aborder de futurs blocs ASIC, avec une meilleure anticipation des risques liés au timing, à la validation et au passage du modèle RTL vers un circuit physique.

\clearpage
